\section{Related Work}
\label{sec:related}

\paragraph{Recruitment extraction resources.} Recruitment datasets differ in their prediction units and annotation targets. Candidate-based classification determines whether a supplied soft-skill phrase refers to an applicant~\citep{sayfullina2018softskills}, while document-level retrieval assigns skills from a predefined inventory to an advertisement~\citep{bhola2020retrieving}. Span-based resources preserve the wording of individual requirements. SkillSpan provides English annotations in potentially overlapping SKILL and KNOWLEDGE layers, covering both hard and soft skills~\citep{zhang2022skillspan}. Kompetencer extends Danish and English skill and knowledge spans with ESCO-derived fine-grained labels~\citep{jensen2022kompetencer}. For German advertisements, \citet{gnehm2022skills} combine requirement-span extraction with finer-grained skill areas and taxonomy matching.

Chinese recruitment research includes entity annotation for occupational profiling and benchmarks for skill-demand forecasting. \citet{cao2021occupational} manually annotated skill and specialty requirements and trained a BERT--BiLSTM--CRF model, while Job-SDF evaluates forecasting at several aggregation levels~\citep{chen2024jobsdf}. Chinese-SkillSpan focuses on extracting source-text requirements as flat character-level spans with four competency types. Table~\ref{tab:related-span} compares selected resources by prediction unit, annotation scheme, and available research materials.

% Compare prediction units, annotation schemes, label sources, and reported materials.
\begin{table}[!htbp]
\centering\small
\caption{Recruitment resources: prediction units, annotation schemes, and materials.}
\label{tab:related-span}
\setlength{\tabcolsep}{4pt}\renewcommand{\arraystretch}{1.13}
\begin{tabularx}{\linewidth}{@{}P{.20\linewidth}P{.17\linewidth}P{.34\linewidth}L@{}}
\toprule
Study / language & Prediction unit & Annotation scheme and label source & Reported materials \\
\midrule
\citet{sayfullina2018softskills}; English & Candidate phrases in context & Binary labels indicating whether a supplied soft-skill phrase refers to an applicant & Annotated soft-skill contexts \\
\citet{gnehm2022skills}; German & Requirement spans and finer-grained skill areas & Coarse EDU/EXP/LNG labels; finer-grained skill areas and taxonomy matching & Article and released domain-adapted models \\
SkillSpan~\citep{zhang2022skillspan}; English & Source-text spans & SKILL and KNOWLEDGE in two potentially overlapping annotation layers & Dataset, code, and annotation guidance \\
Kompetencer~\citep{jensen2022kompetencer}; Danish / English & Spans and fine-grained labels & Manually annotated skill/knowledge spans; distant labels from the ESCO API; manually corrected Danish test labels & Dataset and code \\
\citet{cao2021occupational}; Chinese & Recruitment entities & Manually annotated skill and specialty entities encoded with BIOES & Article and supporting data files \\
Chinese-SkillSpan; Chinese & Flat character-level spans & L/K/S/T types; human reference labels and model-generated silver supervision & Core dataset, handbook, protocols, baseline predictions, scoring code, and individual blinded-annotation layers \\
\bottomrule
\end{tabularx}
\par\smallskip
{\footnotesize\RaggedRight\textit{Notes.} EDU, EXP, and LNG denote education, experience, and language requirements, respectively.\par}
\end{table}


\paragraph{Competency annotation and guidelines.} Competency annotation depends on how boundaries and types are defined in context. Chinese named entity recognition studies address social-media text~\citep{peng2015weibo}, fine-grained categories~\citep{cluener2020}, and boundary detection without explicit word delimiters~\citep{chen2021boundary}. Multilingual benchmarks also examine noisy text and context-dependent entity types~\citep{multiconer2023,dynamicner2025}. In recruitment text, coordinated and ambiguous mentions create annotation difficulties across languages~\citep{nguyen2024rethinking}, while variation in terminology and limited public annotations complicate comparisons among resources~\citep{senger2024survey}.

ESCO provides a multilingual vocabulary for labor-market information~\citep{levrang2014esco,ESCOv1_2_2024}, and SkiLLens combines human involvement with skill extraction and ESCO linking for European job advertisements~\citep{desanto2026skilllens}. Chinese-SkillSpan uses ESCO's high-level distinctions to guide Chinese span annotation. Its annotation rules specify both span boundaries and category assignment~\citep{artstein2008agreement}.

\paragraph{Model-assisted annotation and extraction.} Weak supervision and model assistance offer ways to extend manually annotated resources. Taxonomy-based methods derive labels from skill inventories~\citep{zhang2022weak}, while large language models (LLMs) support automated and collaborative annotation~\citep{ding2023annotator,li2023coannotating}. \citet{ulhaq2026nerannotation} evaluate LLM-generated labels and their utility for downstream model training, including on SkillSpan. GoLLIE uses annotation guidelines for information extraction~\citep{gollie2024}, and \citet{kim2026guidelines} study guideline refinement for LLM annotation in biomedical NER. \mbox{JobSkape} generates synthetic job advertisements for skill-to-taxonomy matching~\citep{magron2024jobskape}. Our silver labels for collected recruitment texts undergo structural checks and sampled human review.

Extraction models draw on both general language representations and recruitment-specific information. Pretrained encoders and structured decoding provide a basis for sequence labeling~\citep{devlin2019bert,lafferty2001crf}, while continued pretraining adapts representations to a target domain~\citep{gururangan2020dapt}. SkillSpan applies this approach to recruitment text through JobBERT and JobSpanBERT~\citep{zhang2022skillspan}. ESCOXLM-R incorporates taxonomy information during multilingual pretraining~\citep{zhang2023escoxlmr}, and NNOSE uses retrieval over training-data representations to support skill extraction~\citep{zhang2024nnose}.

General information extraction methods offer additional approaches to span prediction. UIE uses text-to-structure generation~\citep{uie2022}, UniversalNER learns from distilled supervision~\citep{universalner}, and GLiNER supports open-type extraction with an encoder~\citep{gliner-naacl24}. GPT-NER formulates named entity recognition as generative extraction with self-verification~\citep{gptner2025}. We compare prompted extraction, silver-supervised adaptation, and encoder-based prediction for Chinese competency spans.

\FloatBarrier
